% Traducción española de la transcripción diplomática francesa de EGA IV, cuarta parte, pp. 345--361.
\clearpage
\oldpage[IV]{345}
\phantomsection
\addcontentsline{toc}{section}{Erratas y adendas (Lista 3)}
\section*{ERRATAS Y ADENDAS}
\markboth{ERRATAS Y ADENDAS, LISTA 3}{ERRATAS Y ADENDAS, LISTA 3}
\begin{center}\large (Lista 3)\end{center}

\subsection*{A) Errores tipográficos}

\paragraph{(II, 3.3.2).}
Línea 1 de la p.~58, sustituir $\sh{S}$ por $\sh{M}$. Línea 6 de la
p.~58, sustituir $S(n)$ por $M(n)$.

\paragraph{(II, 4.1.1).}
Línea 8 de la p.~71, sustituir $\sh{O}$-$P$-Module por
$\sh{O}_P$-Module.

\paragraph{(II, 5.4.5).}
Línea 6 de la p.~102, sustituir $u$ por $f\circ u$.

\paragraph{(II, 5.4.8).}
Antes de la línea 4 desde abajo de la p.~102, sustituir $f\times1_{S'}$ por
$f_{(Y')}$ en el diagrama. Línea 4 desde abajo de la p.~102, sustituir
$f\times1_{S'}$ por $f_{(Y')}$.

\paragraph{(II, 6.4.4).}
Línea 1 desde arriba de la p.~122, sustituir normal por intégralement clos.

\paragraph{(II, 6.4.7).}
Línea 14 desde abajo de la p.~122, sustituir dimension por rang.

\paragraph{(II, 6.5.2).}
Línea 5 de la p.~127, sustituir $(\sh{B}|U_\lambda)$ por
$(B|U_\lambda)$. Línea 15 de la p.~127, suprimir el primer
$(N_{\sh{B}/\sh{A}}\omega_{\lambda\mu})$.

\paragraph{($0_{\mathrm{III}}$, 8.1.11).}
Línea 8 desde abajo de la p.~8, sustituir $h_X(u)$ por $h'_X(u)$.

\paragraph{($0_{\mathrm{III}}$, 9.1.2).}
Línea 10 desde abajo de la p.~12, sustituir $\sh{F}$ por $X$.

\paragraph{($0_{\mathrm{III}}$, 10.2.1).}
Línea 6 desde abajo de la p.~18, sustituir « en outre si » por « si en
outre ».

\paragraph{(III, 4.4.12).}
Línea 10 de la p.~138, sustituir chap.~V por chap.~IV (8.11.2).

\paragraph{(III, 5.2.5).}
Línea 12 de la p.~153, sustituir (I, 5.5.4) por (II, 5.5.4).

\paragraph{(III, 7.4.7).}
Línea 15 de la p.~57, sustituir tout $A$-module por tout
$A$-module $M$.

\paragraph{(III, 7.8.5).}
Línea 18 desde abajo de la p.~74, sustituir $F$ por $\sh{F}$.

\paragraph{(III, 7.8.9).}
Línea 9 desde abajo de la p.~75, sustituir $M$ por $\sh{M}$.

\paragraph{(III, 7.9.11).}
Línea 17 desde abajo de la p.~79, sustituir chap.~V por chap.~VI.

\paragraph{($0_{\mathrm{IV}}$, 16.2.2.2).}
Línea 13 de la p.~26, sustituir $M'\cap\mathfrak q^nM'$ por
$M'\cap\mathfrak q^nM$.

\paragraph{($0_{\mathrm{IV}}$, 17.3.7).}
Línea 16 desde abajo de la p.~50, sustituir (17.1.10) por (17.3.3).

\paragraph{($0_{\mathrm{IV}}$, 18.1.3).}
En el diagrama (18.1.3.1), $p_2$ debe estar junto a la flecha
$G\to F$.

\paragraph{($0_{\mathrm{IV}}$, 19.8.7).}
Líneas 7, 8 y 9 desde abajo de la p.~111, sustituir (21.5.5) por (21.5.3).

\paragraph{($0_{\mathrm{IV}}$, 19.8.8).}
Línea 2 desde abajo de la p.~112, sustituir $A$-module por $B$-module.

\paragraph{($0_{\mathrm{IV}}$, 20.5.4).}
Línea 5 desde abajo de la p.~130, sustituir $\Omega_{B'/A'}$ por
$\Omega^1_{B'/A'}$.

\paragraph{($0_{\mathrm{IV}}$, 20.7.8).}
Línea 9 desde abajo de la p.~149, sustituir (18.4.3) por (18.5.4). Línea 4
desde abajo, sustituir (18.4.1) por (18.5.1).

\paragraph{($0_{\mathrm{IV}}$, 21.5.3).}
Línea 1 desde abajo de la p.~166, sustituir ramifiée por ramifié.

\paragraph{($0_{\mathrm{IV}}$, 21.7.1).}
Línea 4 de la p.~170, sustituir (20.6.15.1) por (20.6.17.1).

\paragraph{($0_{\mathrm{IV}}$, 22.1.1).}
Línea 12 de la p.~183, sustituir (20.6.20) por (20.6.10).

\oldpage[IV]{346}
\paragraph{($0_{\mathrm{IV}}$, 22.3.3).}
Línea 13 desde abajo de la p.~192, sustituir $\widehat B_{\mathfrak p}$ por
$\widehat B_{\mathfrak q}$.

\paragraph{($0_{\mathrm{IV}}$, 22.7.3).}
Línea 21 desde abajo de la p.~212, sustituir $\operatorname{Ker}(i)$ por
$\operatorname{Ker}(i_0)$. Línea 12 desde abajo de la p.~212, sustituir
$(\mathfrak p/\mathfrak p^2)\otimes_L L$ por
$(\mathfrak p/\mathfrak p^2)\otimes_A L$.

\paragraph{($0_{\mathrm{IV}}$, 23.1.3).}
Línea 4 de la p.~215, sustituir $x'^q\in K'$ por $x'\in K'$. Línea 9
de la p.~215, sustituir $x'^q\in A'$ por $x'^q\in A$.

\paragraph{({\texorpdfstring{\protect\EGAReadLink{IV.2.1.1}{IV, 2.1.1}}{IV, 2.1.1}}).}
Línea 4 desde abajo de la p.~5, sustituir (III, 1.4.15.1) por
(III, 1.4.15.5).

\paragraph{({\texorpdfstring{\protect\EGAReadLink{IV.2.4.3}{IV, 2.4.3}}{IV, 2.4.3}}).}
Línea 12 de la p.~20, sustituir
$(X_{\mathrm{red}}\times_{Y_{\mathrm{red}}}Y_{\mathrm{red}})_{\mathrm{red}}$
por
$(X_{\mathrm{red}}\times_{Y_{\mathrm{red}}}Y'_{\mathrm{red}})_{\mathrm{red}}$.

\paragraph{({\texorpdfstring{\protect\EGAReadLink{IV.2.6.1}{IV, 2.6.1}}{IV, 2.6.1}}).}
Línea 10 de la p.~27, sustituir $X'\to X_{(S')}$ por
$X'=X_{(S')}$. 

\paragraph{({\texorpdfstring{\protect\EGAReadLink{IV.4.2.1}{IV, 4.2.1}}{IV, 4.2.1}}).}
Línea 20 desde abajo de la p.~55, sustituir
$\operatorname{deg.tr.}_kE$ por $\operatorname{deg.tr.}_LE$ y
$\operatorname{deg.tr.}_kL(t)$ por $\operatorname{deg.tr.}_LL(t)$;
línea 19 desde abajo, intercambiar « première » y « seconde ».

\paragraph{({\texorpdfstring{\protect\EGAReadLink{IV.4.4.2}{IV, 4.4.2}}{IV, 4.4.2}}).}
Línea 1 desde abajo de la p.~59, sustituir $X$ por $X'$.

\paragraph{({\texorpdfstring{\protect\EGAReadLink{IV.4.7.3}{IV, 4.7.3}}{IV, 4.7.3}}).}
Línea 10 desde abajo de la p.~75, sustituir préschéma por $k$-préschéma.

\paragraph{({\texorpdfstring{\protect\EGAReadLink{IV.5.6.1}{IV, 5.6.1}}{IV, 5.6.1}}).}
Línea 16 de la p.~98, sustituir (en el subíndice) $k(\mathfrak p($ por
$k(\mathfrak p)$.

\paragraph{({\texorpdfstring{\protect\EGAReadLink{IV.5.6.5}{IV, 5.6.5}}{IV, 5.6.5}}).}
Línea 20 de la p.~99, suprimir el segundo paréntesis en $k(\xi))$.

\paragraph{({\texorpdfstring{\protect\EGAReadLink{IV.5.6.11.1}{IV, 5.6.11.1}}{IV, 5.6.11.1}}).}
Línea 16 de la p.~102, sustituir $K$ por $\mathfrak K$.

\paragraph{({\texorpdfstring{\protect\EGAReadLink{IV.5.11.2}{IV, 5.11.2}}{IV, 5.11.2}}).}
Nota al pie de la p.~123, sustituir 5.11.2) por (5.11.2).

\paragraph{({\texorpdfstring{\protect\EGAReadLink{IV.6.1.5}{IV, 6.1.5}}{IV, 6.1.5}}).}
Línea 2 desde abajo de la p.~136, sustituir (5.5.1.2) por
($0$, 16.3.9.2).

\paragraph{({\texorpdfstring{\protect\EGAReadLink{IV.6.3.1}{IV, 6.3.1}}{IV, 6.3.1}}).}
Línea 9 de la p.~139, sustituir $y\in\mathfrak m$ por
$y\in\mathfrak n$.

\paragraph{({\texorpdfstring{\protect\EGAReadLink{IV.6.7.2}{IV, 6.7.2}}{IV, 6.7.2}}).}
Línea 18 de la p.~146, sustituir $\sh{O}'_{x'}$ por $\sh{O}_{x'}$.

\paragraph{({\texorpdfstring{\protect\EGAReadLink{IV.6.10.1}{IV, 6.10.1}}{IV, 6.10.1}}).}
Línea 13 de la p.~155, sustituir
$\sh{F}\otimes_{\sh{O}_Y}\sh{O}_Y$ por
$\sh{F}\otimes_{\sh{O}_X}\sh{O}_Y$.

\paragraph{({\texorpdfstring{\protect\EGAReadLink{IV.6.14.1}{IV, 6.14.1}}{IV, 6.14.1}}).}
Línea 21 de la p.~170, sustituir n\textsuperscript{o}~3, prop.~12
por n\textsuperscript{o}~2, prop.~8.

\paragraph{({\texorpdfstring{\protect\EGAReadLink{IV.6.15.2}{IV, 6.15.2}}{IV, 6.15.2}}).}
Línea 13 de la p.~177, sustituir $t(t-v)$ por $t^2(t-v)^2$.

\paragraph{({\texorpdfstring{\protect\EGAReadLink{IV.6.15.7}{IV, 6.15.7}}{IV, 6.15.7}}).}
Línea 6 desde abajo de la p.~179, sustituir $X\otimes_Ak'$ por
$X\otimes_kk'$.

\paragraph{({\texorpdfstring{\protect\EGAReadLink{IV.7.5.4}{IV, 7.5.4}}{IV, 7.5.4}}).}
Línea 14 desde abajo de la p.~206, sustituir
$\operatorname{Spec}(\widehat B\otimes k)$ por
$\operatorname{Spec}(\widehat B\otimes_{\widehat A}k)$.

\paragraph{({\texorpdfstring{\protect\EGAReadLink{IV.7.8.4}{IV, 7.8.4}}{IV, 7.8.4}}).}
Línea 10 desde abajo de la p.~216, sustituir « vérifiant »
({\EGAReadLink{IV.7.8.2}{7.8.2}}, (ii) y (iii)) por
« vérifiant ({\EGAReadLink{IV.7.8.2}{7.8.2}}, (ii) et (iii)) ».

\paragraph{(IV, 2\textsuperscript{e} Partie, Bibliographie).}
Línea 6 de la p.~224, sustituir \emph{Amer. J. of Math.} por
\emph{Annals of Math.}

\paragraph{({\texorpdfstring{\protect\EGAReadLink{IV.8.5.8}{IV, 8.5.8}}{IV, 8.5.8}}).}
Línea 6 de la p.~24, sustituir $\operatorname{Ker}g_\mu$ al final de
la línea por $\operatorname{Ker}g_\lambda$.

\paragraph{({\texorpdfstring{\protect\EGAReadLink{IV.8.10.1}{IV, 8.10.1}}{IV, 8.10.1}}).}
Línea 13 desde abajo de la p.~36, sustituir (2.3.11) por
{\EGAReadLink{IV.8.2.11}{(8.2.11)}}.

\paragraph{({\texorpdfstring{\protect\EGAReadLink{IV.10.2.1}{IV, 10.2.1}}{IV, 10.2.1}}).}
Línea 14 desde abajo de la p.~98, sustituir $y$ por $Y$.

\paragraph{({\texorpdfstring{\protect\EGAReadLink{IV.10.4.11}{IV, 10.4.11}}{IV, 10.4.11}}).}
Línea 10 desde abajo de la p.~104, sustituir (17.9.7) por
{\EGAReadLink{IV.17.9.6}{(17.9.6)}}.

\paragraph{({\texorpdfstring{\protect\EGAReadLink{IV.11.2.9}{IV, 11.2.9}}{IV, 11.2.9}}).}
Línea 3 desde abajo de la p.~126, sustituir $A_\lambda$ por $A_\mu$.

\paragraph{({\texorpdfstring{\protect\EGAReadLink{IV.11.3.4}{IV, 11.3.4}}{IV, 11.3.4}}).}
Línea 16 de la p.~134, sustituir $F$ por $\sh{F}$.

\paragraph{({\texorpdfstring{\protect\EGAReadLink{IV.11.3.10}{IV, 11.3.10}}{IV, 11.3.10}}).}
Línea 9 desde abajo de la p.~139, sustituir (11.3.7) por
{\EGAReadLink{IV.11.3.10}{(11.3.10)}}.

\paragraph{({\texorpdfstring{\protect\EGAReadLink{IV.11.7.5}{IV, 11.7.5}}{IV, 11.7.5}}).}
Línea 11 de la p.~159, sustituir (6.5.11) por (6.15.11).

\paragraph{({\texorpdfstring{\protect\EGAReadLink{IV.12.1.1}{IV, 12.1.1}}{IV, 12.1.1}}).}
Línea 13 de la p.~174, añadir al final: pour $k\geq1$.

\paragraph{({\texorpdfstring{\protect\EGAReadLink{IV.14.2.1}{IV, 14.2.1}}{IV, 14.2.1}}).}
Línea 11 de la p.~202, sustituir (6.1.2) por (6.1.1).

\oldpage[IV]{347}
\paragraph{({\texorpdfstring{\protect\EGAReadLink{IV.14.3.8}{IV, 14.3.8}}{IV, 14.3.8}}).}
Línea 10 desde abajo de la p.~206, sustituir $(f^{-1}y)$ por $f^{-1}(y)$.

\paragraph{({\texorpdfstring{\protect\EGAReadLink{IV.14.5.7}{IV, 14.5.7}}{IV, 14.5.7}}).}
Línea 16 de la p.~218, sustituir (14.5.5) por (14.5.6).

\paragraph{({\texorpdfstring{\protect\EGAReadLink{IV.14.5.11}{IV, 14.5.11}}{IV, 14.5.11}}).}
En el diagrama de la parte superior de la página 223, todas las flechas
horizontales deben ir de derecha a izquierda.

\paragraph{({\texorpdfstring{\protect\EGAReadLink{IV.15.4.2}{IV, 15.4.2}}{IV, 15.4.2}}).}
Línea 23 de la p.~230, sustituir (6.1.4) por (6.1.5).

\paragraph{({\texorpdfstring{\protect\EGAReadLink{IV.15.5.2}{IV, 15.5.2}}{IV, 15.5.2}}).}
Línea 16 de la p.~233, sustituir $f^{-1}(y)$ por $f^{-1}(y')$.

\paragraph{({\texorpdfstring{\protect\EGAReadLink{IV.15.6.10}{IV, 15.6.10}}{IV, 15.6.10}}).}
Línea 6 de la p.~242, sustituir ${}^{-1}(y')$ por $f^{-1}(y')$.
Línea 11 de la p.~242, sustituir $\alpha$) et $\beta$) por
$\alpha$) ou $\beta$).

\subsection*{B) Modificaciones del texto}

\paragraph{(Err\textsubscript{IV}, 1).}
En ($0_{\mathrm{III}}$, 11.1.6), el enunciado (que reproduce un enunciado de
(G, I, 4.4.1)) es incorrecto cuando $q(n)$ es un entero arbitrario;
solo es correcto si $q(n)=0$ o $q(n)=n$. En efecto, en el
primer caso, el núcleo de $d_r^{p,0}$ es entonces $E_r^{p,0}$, pues
la imagen de $d_r^{p,0}$ es $E_r^{p+r,-r+1}=0$ por hipótesis; asimismo,
la imagen de $d_r^{p-r,r-1}$ es nula puesto que
$E_r^{p-r,r-1}=0$; se tiene, por tanto, $E_{r+1}^{p,0}=E_r^{p,0}$. Se
razona del mismo modo cuando $q(n)=n$ para todo $n$.

\paragraph{(Err\textsubscript{IV}, 2).}
En (III, 3.4), suprimir las líneas 2 y 3 de la p.~119.

\oldpage[IV]{348}
\paragraph{(Err\textsubscript{IV}, 3).}
En (III, 4.4.9), línea 15 desde abajo, p.~137, sustituir finies por
finies en tant qu'ensembles.

\paragraph{(Err\textsubscript{IV}, 4).}
En (III, 7.8.10), línea 9 de la p.~76, sustituir simple por lisse.

\paragraph{(Err\textsubscript{IV}, 5).}
En ($0_{\mathrm{IV}}$, 16.1.4), línea 8 desde abajo de la p.~23, sustituir
Dire que $A$ est caténaire por Dire qu'un anneau intègre $A$ est
caténaire.

\paragraph{(Err\textsubscript{IV}, 6).}
En ($0_{\mathrm{IV}}$, 16.1.9), puede suprimirse la hipótesis de que
$A$ sea noetheriano. En efecto, para todo elemento $u\in B_M$, $A[u]$
es una sub-$A$-álgebra de $C$, luego $u$ es entero sobre $A$ (Bourbaki,
\emph{Alg. comm.}, cap.~V, \S1, n\textsuperscript{o}~1, cor. de la
prop.~1), y por consiguiente también sobre $A_M$; se concluye igualmente mediante
{\EGAReadLink{IV.16.1.5}{(16.1.5)}}.

\paragraph{(Err\textsubscript{IV}, 7).}
En ($0_{\mathrm{IV}}$, 16.5.12), obsérvese que la relación
(16.5.11.1) equivale a la misma relación en la que se sustituye $A$ por los
anillos cociente $A/\mathfrak r$, al recorrer $\mathfrak r$ el conjunto
de los ideales primos minimales de $A$. Por tanto, en la demostración de
{\EGAReadLink{IV.16.5.12}{(16.5.12)}} podemos limitarnos al caso
en que $A$ sea íntegro y aplicar (16.1.4.2).

\paragraph{(Err\textsubscript{IV}, 8).}
Suprimir ($0_{\mathrm{IV}}$, 19.3.12), cuya demostración es errónea
(los conjuntos $W_\lambda$, $U_\lambda$ y $V_\lambda$ no son
variedades lineales).

\paragraph{(Err\textsubscript{IV}, 9).}
En ($0_{\mathrm{IV}}$, 19.8.3), líneas 6 y 8 desde abajo de la p.~110,
sustituir homomorphisme por homomorphisme local; lo mismo en
($0_{\mathrm{IV}}$, 19.8.4), línea 1 de la p.~111.

\paragraph{(Err\textsubscript{IV}, 10).}
En ($0_{\mathrm{IV}}$, 20.1.2), línea 2 desde abajo de la p.~117,
sustituir $a\in A$ por $b\in B$ y $aD$ por $bD$; línea 1 desde abajo de la
p.~117, sustituir $A$-module por $B$-module.

\paragraph{(Err\textsubscript{IV}, 11).}
En ($0_{\mathrm{IV}}$, 20.2.1), línea 8 de la p.~119, sustituir
$A$-modules por $B$-modules. Líneas 21 y 22 de la p.~119, sustituir
« di-homomorphismes de modules (relatif à $u$) » por homomorphisme
de $B$-modules.

\paragraph{(Err\textsubscript{IV}, 12).}
En ($0_{\mathrm{IV}}$, 20.4.6.1), sustituir $p_2-p_1$ por $p_1-p_2$;
línea 10 desde abajo de la p.~125, sustituir $x\otimes1-1\otimes x$ por
$1\otimes x-x\otimes1$.

\paragraph{(Err\textsubscript{IV}, 13).}
En ($0_{\mathrm{IV}}$, 20.4.8), línea 14 de la p.~126, sustituir
$p_2=p_1-d_{B/A}$ por $p_1=p_2-d_{B/A}$.

\paragraph{(Err\textsubscript{IV}, 14).}
En ($0_{\mathrm{IV}}$, 20.4.13), sustituir las líneas 6 a 11 desde abajo
de la p.~127 por el texto siguiente:

\medskip
\noindent
(iv) Sean $A$, $B$ dos anillos íntegros tales que $A\subset B$, siendo $B$
entero sobre $A$, y que el cuerpo de fracciones de $B$ sea una
extensión separable del de $A$. Entonces el $B$-módulo
$\Omega^1_{B/A}$ es un módulo de torsión. En efecto, para todo $x\in B$,
sea $f$ el polinomio minimal de $x$ respecto del cuerpo de fracciones
$K$ de $A$; existe un $a\in A$ no nulo tal que $af[T]\in A[T]$; como
$x$ es separable sobre $K$, se tiene $f'(x)\neq0$, y, por otra parte, de la
relación $af(x)=0$ se deduce $af'(x)d_{B/A}x=0$, de donde se sigue nuestra afirmación
en virtud de {\EGAReadLink{IV.20.4.7}{(20.4.7)}}.

\paragraph{(Err\textsubscript{IV}, 15).}
En ($0_{\mathrm{IV}}$, 22.2.7), línea 5 de la p.~189, en lugar de
« formellement lisses », debe decir « formellement lisses pour les
topologies préadiques ».

\paragraph{(Err\textsubscript{IV}, 16).}
En el Sumario del cap.~IV, línea 4 desde abajo de la p.~222, añadir:
et anneaux strictement locaux. Sustituir las líneas 1 a 3 desde abajo por:
\begin{quote}
\S19. Inmersiones regulares y planitud normal.\\
\S20. Funciones meromorfas y seudomorfismos.\\
\S21. Divisores.
\end{quote}
Línea 16 de la p.~223, sustituir « de présentation finie » por
« localement de présentation finie »; línea 19 de la p.~223, sustituir
« comme » por « comme étant ».

\paragraph{(Err\textsubscript{IV}, 17).}
En (IV, 3.1.10), añadir después de la línea 16 de la p.~39:

\begin{remark}[3.1.10.2]
\label{IV.3.1.10.2}
Sea $f:X\to Y$ un morfismo cuasicompacto y cuasiseparado, y sea
$\sh{F}$ un $\sh{O}_X$-módulo cuasicoherente, de modo que
$f_*(\sh{F})$ es un $\sh{O}_Y$-módulo cuasicoherente
{\EGAReadLink{IV.1.7.4}{(1.7.4)}}. Entonces se tiene
\begin{align}
 \operatorname{Ass}(f_*(\sh{F}))&\subset f(X),
 \label{IV.3.1.10.3}\tag{3.1.10.3}\\
 \operatorname{Ass}(f_*(\sh{F}))&\subset f(\operatorname{Ass}(\sh{F}))
 \quad\text{si $X$ es localmente noetheriano}.
 \label{IV.3.1.10.4}\tag{3.1.10.4}
\end{align}

Sea, en efecto, $y\in\operatorname{Ass}(f_*(\sh{F}))$; en virtud de
{\EGAReadLink{IV.2.3.1}{(2.3.1)}} y del hecho de que el morfismo
$\operatorname{Spec}(\sh{O}_{Y,y})\to Y$ es plano,
$(f_*(\sh{F}))_y$ es igual a $(f'_*(\sh{F}'))_y$, donde
\[
 f':X\times_Y\operatorname{Spec}(\sh{O}_{Y,y})
 \longrightarrow\operatorname{Spec}(\sh{O}_{Y,y})
\]
se deduce de $f$ por cambio de base y
$\sh{F}'=\sh{F}\otimes_{\sh{O}_Y}\sh{O}_{Y,y}$. Quedamos, por tanto, reducidos a
demostrar (3.1.10.3) y (3.1.10.4), sustituyendo en ellas $Y$ por
$\operatorname{Spec}(\sh{O}_{Y,y})$, $f$ por $f'$ y $\sh{F}$ por
$\sh{F}'$, teniendo en cuenta (I, 3.6.5) y la definición de
$\operatorname{Ass}(\sh{F}')$.

Por definición, existe un elemento
$t_y\in(f_*(\sh{F}))_y$ no nulo cuyo anulador es $\mathfrak m_y$, y como
$(f_*(\sh{F}))_y$ se identifica con
$\Gamma(Y,f_*(\sh{F}))$ puesto que $Y$ es local, $t_y$ se identifica asimismo
con una sección $s\in\Gamma(X,\sh{F})$. El ideal de $\sh{O}_X$ anulador
de $s$ contiene, por tanto, $\mathfrak m_y\sh{O}_X$, y se deduce de
($0_{\mathrm I}$, 1.7.4) que el soporte de $s$ está contenido en
$f^{-1}(y)$; si este último fuese vacío, se tendría $s=0$, lo cual es
absurdo. Si además $X$ es localmente noetheriano, el
sub-$\sh{O}_X$-módulo coherente $\sh{O}_Xs$ de $\sh{F}$ no es nulo,
luego $\operatorname{Ass}(\sh{O}_Xs)\neq\varnothing$
{\EGAReadLink{IV.3.1.5}{(3.1.5)}}; pero
\[
 \operatorname{Ass}(\sh{O}_Xs)\subset
 \operatorname{Supp}(\sh{O}_Xs)\subset f^{-1}(y),\qquad
 \operatorname{Ass}(\sh{O}_Xs)\subset\operatorname{Ass}(\sh{F})
\]
{\EGAReadLink{IV.3.1.7}{(3.1.7)}}, por tanto
$y\in f(\operatorname{Ass}(\sh{F}))$.
\end{remark}

\oldpage[IV]{349}
\begin{remark}[3.1.10.5]
\label{IV.3.1.10.5}
Obsérvese que, aun cuando $f$ es un morfismo propio, no se tiene
necesariamente
\[
 f(\operatorname{Ass}(\sh{F}))\subset
 \operatorname{Ass}(f_*(\sh{F})).
\]
Esto se ve tomando como $Y$ el espectro de un cuerpo y como $X$ un
espacio proyectivo sobre $k$, por ejemplo $\mathbf P^1_k$; entonces
$\sh{F}=\sh{O}_X(-1)$ es un $\sh{O}_X$-módulo coherente $\neq0$ tal
que $\Gamma(X,\sh{F})=0$ (III, 2.1.13), luego
$\operatorname{Ass}(f_*(\sh{F}))=\varnothing$ y
$\operatorname{Ass}(\sh{F})\neq\varnothing$.
\end{remark}

\paragraph{(Err\textsubscript{IV}, 18).}
Después de (IV, 3.2.8), insertar después de la línea 17 de la p.~43:

\begin{remark}[3.2.9]
\label{IV.3.2.9}
Para todo punto $s$ de un preesquema $X$, denotemos por
$i_s:\operatorname{Spec}(k(s))\to X$ el morfismo canónico, que es
un monomorfismo cuasicompacto. Para todo $k(s)$-módulo $G(s)$,
considerado como módulo sobre $\operatorname{Spec}(k(s))$,
$(i_s)_*(G(s))$ es, por tanto, un $\sh{O}_X$-módulo cuasicoherente
{\EGAReadLink{IV.1.7.4}{(1.7.4)}}. Se comprueba inmediatamente
que la fibra de $(i_s)_*(G(s))$ es nula en todo punto
$x\notin\overline{\{s\}}$ e isomorfa a $G(s)$ en caso contrario.
Si $X$ es localmente noetheriano y $G(s)$ es de longitud finita, está
claro que $\operatorname{Ass}((i_s)_*(G(s)))$ se reduce al único punto
$s$; en otras palabras, $(i_s)_*(G(s))$ es irredundante.

Supongamos ahora que $X$ es localmente noetheriano, y sea $\sh{F}$ un
$\sh{O}_X$-módulo coherente. Sea, por otra parte, $S$ una parte discreta
de $X$, y, para todo $s\in S$, pongamos
$\sh{F}(s)=i_s^*(\sh{F})$, que es un $k(s)$-módulo; consideremos para
cada $s\in S$ un módulo cociente $G(s)$ de $\sh{F}(s)$ de longitud
finita. Se tiene entonces un homomorfismo canónico
\[
 u:\sh{F}\longrightarrow
 \bigoplus_{s\in S}(i_s)_*(\sh{F}(s))\longrightarrow
 \bigoplus_{s\in S}(i_s)_*(G(s))
\]
de $\sh{O}_X$-módulos coherentes (pues $S$, al ser discreto, es
localmente finito en $X$). Si $\sh{H}=\operatorname{Ker}u$,
$\sh{G}=\sh{F}/\sh{H}$ es un $\sh{O}_X$-módulo coherente y se deduce
de $u$ un homomorfismo inyectivo
\[
 v:\sh{G}\longrightarrow\bigoplus_{s\in S}(i_s)_*(G(s)).
\]
Si $\sh{G}^{(s)}$ es la imagen de $\sh{G}$ en $(i_s)_*(G(s))$,
$(\sh{G}^{(s)})_{s\in S}$ es una descomposición irredundante, que es
reducida cuando todos los $G(s)$ son $\neq0$, y tal que
$\operatorname{Ass}(\sh{G})=S$ (por tanto, $\sh{G}$ no tiene ciclo primo
asociado sumergido). Se obtiene así una biyección
$(G(s))_{s\in S}\leftrightarrow\sh{G}$, donde $S$ recorre el conjunto de las
partes discretas de $X$ para las cuales $\sh{F}(s)\neq0$ para todo
$s\in S$ y, para cada $S$, $(G(s))_{s\in S}$ recorre el conjunto
de los cocientes $\neq0$ y de longitud finita de los $\sh{F}(s)$; por otra
parte, $\sh{G}$ recorre el conjunto de los $\sh{O}_X$-módulos cocientes
coherentes de $\sh{F}$, sin ciclo primo sumergido. La biyección
recíproca asocia a cada $\sh{G}$ su descomposición irredundante
reducida.
\end{remark}

\paragraph{(Err\textsubscript{IV}, 19).}
En (IV, 4.2.1), antes de la línea 17 desde abajo de la p.~55, insertar:

\begin{remark}[4.2.1.4]
\label{IV.4.2.1.4}
Bajo las hipótesis de {\EGAReadLink{IV.4.2.1}{(4.2.1)}}, se tiene
además
\begin{equation}
 \dim(K\otimes_kL)=
 \inf(\operatorname{deg.tr.}_kK,\operatorname{deg.tr.}_kL).
 \label{IV.4.2.1.5}\tag{4.2.1.5}
\end{equation}
donde debe entenderse que, si los dos cuerpos $K$, $L$ son de grado de
trascendencia infinito (que pueden ser cardinales arbitrarios), el
primer miembro debe tomarse igual a $+\infty$ (la dimensión considerada
aquí es la definida en ($0$, 16.1.1)).

Supongamos, en efecto, por ejemplo, que
$\operatorname{deg.tr.}_kK\leq\operatorname{deg.tr.}_kL$ y, en primer
lugar, que $\operatorname{deg.tr.}_kK=n<+\infty$. Se ha visto en la demostración
de {\EGAReadLink{IV.4.2.1}{(4.2.1)}} que $K\otimes_kL$ es una
extensión algebraica de $K'\otimes_kL$ y tiene, por tanto, la misma dimensión
($0$, 16.1.5). Podemos, pues, limitarnos al caso en que
$K=k(T_1,\ldots,T_n)$\footnote{La lista de erratas impresa y el inglés
mantenido usan corchetes en esta definición; la lectura corregida usa
paréntesis porque $K$ es el cuerpo de funciones racionales.}; entonces
$K\otimes_kL$ es un anillo de fracciones
de
\[
 k[T_1,\ldots,T_n]\otimes_kL=L[T_1,\ldots,T_n],
\]
que es de dimensión $n$ {\EGAReadLink{IV.5.5.4}{(5.5.4)}}, y
por consiguiente $\dim(K\otimes_kL)\leq n$ ($0$, 16.1.3.1). Para demostrar que
$\dim(K\otimes_kL)=n$, basta evidentemente mostrar que existe en
$\operatorname{Spec}(L[T_1,\ldots,T_n])$ un punto racional sobre $L$
que pertenece a $\operatorname{Spec}(K\otimes_kL)$
{\EGAReadLink{IV.5.2.3}{(5.2.3)}}, o, equivalentemente (I, 1.7.4), que
existe un $L$-homomorfismo $u:K\otimes_kL\to L$. Pero, por definición
del producto tensorial, esto equivale a demostrar la existencia de un
$k$-homomorfismo $v:K\to L$, lo cual es evidente en virtud de
la hipótesis
$\operatorname{deg.tr.}_kK\leq\operatorname{deg.tr.}_kL$.

Si $\operatorname{deg.tr.}_kK\leq\operatorname{deg.tr.}_kL$ y estos
dos cardinales son infinitos, se ve del mismo modo que existe un punto
racional en $\operatorname{Spec}(L[(t_\alpha)_{\alpha\in I}])$ que
pertenece a $\operatorname{Spec}(K\otimes_kL)$; pero el núcleo de todo
$L$-homomorfismo de $L[(t_\alpha)_{\alpha\in I}]$ en $L$ es un
ideal maximal engendrado por una familia de la forma
$(t_\alpha-a_\alpha)_{\alpha\in I}$ con $a_\alpha\in L$. Está claro
que tal ideal maximal contiene una sucesión infinita decreciente
de ideales primos distintos, luego el anillo local correspondiente de
$L[(t_\alpha)_{\alpha\in I}]$ no es noetheriano y tiene
dimensión infinita; como es un anillo local de $K\otimes_kL$, se
ve que este último es no noetheriano y de dimensión infinita.
\end{remark}

\oldpage[IV]{350}
\paragraph{(Err\textsubscript{IV}, 20).}
Después de (IV, 4.5.19), insertar, antes de la línea 12 desde abajo de la p.~67:

\begin{corollary}[4.5.19.3]
\label{IV.4.5.19.3}
Sean $k$ un cuerpo y $X$ un $k$-preesquema de tipo finito sobre $k$, que sea
irreducible pero no geométricamente irreducible. Entonces existe una
parte cerrada $Y$ de $X$, rara en $X$, que contiene todos los puntos
$x\in X$ tales que $k(x)$ sea radicial sobre $k$ (y, en particular, todos
los puntos de $X$ racionales sobre $k$).
\end{corollary}

Sean $k'$ una clausura algebraica de $k$, $X'=X\otimes_kk'$,
$p:X'\to X$ la proyección canónica. Si $x\in X$ es tal que $k(x)$
sea radicial sobre $k$, solo existe un punto $x'\in X'$ sobre
$x$ (I, 3.4.9); este punto pertenece necesariamente a dos componentes
irreducibles distintas de $X'$, pues de otro modo $X$ sería geométricamente
irreducible en virtud de
{\EGAReadLink{IV.4.5.19}{(4.5.19)}}, en contra de
la hipótesis. Sean entonces $X'_i$ ($1\leq i\leq n$) las componentes
irreducibles del preesquema noetheriano $X'$ y, para
$1\leq i<j\leq n$, pongamos $X'_{ij}=X'_i\cap X'_j$; el conjunto
\[
 Y=p\left(\bigcup_{i<j}X'_{ij}\right)
\]
responde a la cuestión. En efecto, como $p$ es un morfismo entero y
los $X'_{ij}$ son cerrados en $X'$, $Y$ es cerrado en $X$.
Por otra parte, si $\xi$ es el punto genérico de $X$, $p^{-1}(\xi)$
es discreto y está formado por los puntos genéricos de los $X'_i$
{\EGAReadLink{IV.4.5.11}{(4.5.11)}}, por lo que no corta ninguno
de los $X'_{ij}$; por consiguiente, el conjunto cerrado $Y$ es raro.

\paragraph{(Err\textsubscript{IV}, 21).}
En (IV, 5.6.11.1), línea 7 desde abajo de la p.~102, después de « intègre »,
añadir « de dimension 2 ».

\paragraph{(Err\textsubscript{IV}, 22).}
Después de (IV, 6.2.3), después de la línea 18 de la p.~138, insertar:

\begin{proposition}[6.2.4]
\label{IV.6.2.4}
Bajo las hipótesis de {\EGAReadLink{IV.6.2.3}{(6.2.3)}}, sea
además $M$ un $A$-módulo de tipo finito y supongamos $M\neq0$ y
$N\neq0$. Entonces se tiene
\begin{equation}
 \operatorname{dim.proj}_B(M\otimes_AN)=
 \operatorname{dim.proj}_A(M)+
 \operatorname{dim.proj}_{B\otimes_Ak}(N\otimes_Ak).
 \label{IV.6.2.4.1}\tag{6.2.4.1}
\end{equation}
\end{proposition}

Consideremos la resolución por la izquierda $L_\bullet$ de $N$ introducida en
{\EGAReadLink{IV.6.2.3}{(6.2.3)}} y una resolución por la izquierda
$K_\bullet$ de $M$ por $A$-módulos libres de tipo finito:
\[
 \cdots\longrightarrow K_i\longrightarrow K_{i-1}\longrightarrow
 \cdots\longrightarrow K_0\longrightarrow M\longrightarrow0.
\]
Se ha visto que los $L_i$, $N$, los $Z_i(L_\bullet)$ son $A$-módulos
planos; lo mismo ocurre con los $B_i(L_\bullet)$, que son iguales a los
$Z_i(L_\bullet)$ salvo para $i=0$, y en este último caso
$B_0(L_\bullet)$ es

\oldpage[IV]{351}
plano por ($0_{\mathrm I}$, 6.1.2); por último, los $H_i(L_\bullet)$ son nulos
por definición salvo $H_0(L_\bullet)=N$, que es un $A$-módulo plano. Se
sabe que, en estas condiciones,
\[
 H_n(K_\bullet\otimes_AL_\bullet)\simeq
 \bigoplus_{p+q=n}
 H_p(K_\bullet)\otimes_AH_q(L_\bullet)
\]
(G, I, 5.5), luego $K_\bullet\otimes_AL_\bullet$ (graduado como de
costumbre) es una resolución de $M\otimes_AN$ por $B$-módulos
libres. Pongamos
\[
 r=\operatorname{dim.proj}_A(M),\qquad
 s=\operatorname{dim.proj}_B(N)=
 \operatorname{dim.proj}_{B\otimes_Ak}(N\otimes_Ak).
\]
Existe, por tanto, una resolución libre $K_\bullet$ de longitud $r$ y una
resolución libre $L_\bullet$ de longitud $s$, y
$K_\bullet\otimes_AL_\bullet$ es entonces una resolución libre de longitud
$r+s$, lo cual demuestra ya que
$\operatorname{dim.proj}_B(M\otimes_AN)\leq r+s$. Queda por demostrar,
según ($0$, 17.2.6), que se tiene
\[
 \operatorname{Tor}^{B}_{r+s}(M\otimes_AN,k(B))\neq0.
\]
Utilizando la resolución precedente $K_\bullet\otimes_AL_\bullet$ para
calcular este funtor, se obtiene
\[
 \operatorname{Tor}^{B}_{r+s}(M\otimes_AN,k(B))=
 H_{r+s}((K_\bullet\otimes_AL_\bullet)\otimes_Bk(B)).
\]
Pero puede escribirse
\[
 (K_\bullet\otimes_AL_\bullet)\otimes_Bk(B)=
 (K_\bullet\otimes_Ak)\otimes_k
 (L_\bullet\otimes_Bk(B)),
\]
que es, por tanto, un producto tensorial de dos complejos de espacios
vectoriales sobre $k$; como todos los $k$-módulos son planos, la fórmula de
Künneth da
\[
 H_{r+s}((K_\bullet\otimes_Ak)\otimes_k
 (L_\bullet\otimes_Bk(B)))\simeq
 H_r(K_\bullet\otimes_Ak)\otimes_k
 H_s(L_\bullet\otimes_Bk(B)),
\]
siendo nula la homología del primer (resp. segundo) complejo en grados
$>r$ (resp. $>s$). Se tiene, por tanto,
\[
 \operatorname{Tor}^{B}_{r+s}(M\otimes_AN,k(B))=
 \operatorname{Tor}^{A}_{r}(M,k)\otimes_k
 \operatorname{Tor}^{B}_{s}(N,k(B))
\]
salvo isomorfismo, y como cada uno de los factores es $\neq0$ por
hipótesis, lo mismo ocurre con su producto tensorial sobre $k$.

\paragraph{(Err\textsubscript{IV}, 23).}
En (IV, 6.10.1), línea 8 de la p.~155, después de canonique, añadir:
$\sh{F}$ un $\sh{O}_X$-Module cohérent.

\paragraph{(Err\textsubscript{IV}, 24).}
En ({\EGAReadLink{IV.6.14.4}{IV, 6.14.4}}), la referencia
dada en la línea 12 desde abajo de la p.~175 es insuficiente; hay que razonar
del modo siguiente. En primer lugar, si $M$ es una base de trascendencia
de $K'$ sobre $A$, $A(M)$ es algebraicamente cerrado en $B(A(M))$
(Bourbaki, \emph{Alg.}, cap.~V, \S6, ejerc.~8 c)). Podemos, por tanto,
limitarnos al caso en que $K'$ sea una extensión algebraica separable de $A$.
Sea $p$ el exponente característico de $A$; se deduce de Bourbaki,
\emph{loc. cit.}, \S9, ejerc.~2, que la clausura algebraica de $K'$
en $B(K')$ es radicial sobre $K'$; demostremos que todo elemento
$x\in B(K')$ tal que $x^{p^r}\in K'$ pertenece ya a $K'$. Para ello,
consideremos una base $(e_\lambda)$ de $K'$ sobre $A$; puesto que $K'$ es una
extensión algebraica separable de $A$, los $e_\lambda^{p^r}$ forman
también una base de $K'$ sobre $A$ (Bourbaki, \emph{loc. cit.}, \S8,
n\textsuperscript{o}~3, cor. de la prop.~4). Se tiene
$x=\sum_\lambda e_\lambda b_\lambda$, con $b_\lambda\in B$ y, por
hipótesis, existen elementos $a_\lambda\in A$ tales que
\[
 \sum_\lambda e_\lambda^{p^r}b_\lambda^{p^r}=
 \sum_\lambda e_\lambda^{p^r}a_\lambda;
\]
pero esto implica $b_\lambda^{p^r}=a_\lambda$ para todo $\lambda$;
como por hipótesis $A$ es algebraicamente cerrado en $B$, se tiene, por tanto,
$b_\lambda\in A$ para todo $\lambda$ y, por consiguiente, $x\in K'$.

\paragraph{(Err\textsubscript{IV}, 25).}
En (IV, 7.4.8, B)), p.~203, la cuestión ha sido resuelta afirmativamente
por R.~Kiehl, \emph{Ausgezeichnete Ringe der nichtarchimedischen
analytischen Geometrie}, de próxima aparición en \emph{Inventiones
Mathematicae}.

\paragraph{(Err\textsubscript{IV}, 26).}
En (IV, 7.8.4), línea 9 desde abajo de la p.~217, sustituir:
\begin{quote}
Veremos más adelante (\S18) que si $k$ es un cuerpo valuado completo no
discreto\ldots
\end{quote}
por:
\begin{quote}
R.~Kiehl ha demostrado (en el artículo citado en
(Err\textsubscript{IV}, 25)) que si $k$ es un cuerpo valuado completo no
discreto\ldots
\end{quote}

\oldpage[IV]{352}
\paragraph{(Err\textsubscript{IV}, 27).}
En (IV, 7.8.4), añadir, antes de la línea 6 desde abajo de la p.~217:

\smallskip
\noindent
\textup{(vi)} Sea $A$ un anillo local noetheriano íntegro de dimensión
1; es entonces universalmente catenario
{\EGAReadLink{IV.7.2.9}{(7.2.9)}} y, para que sea excelente,
es necesario y suficiente, en virtud de
{\EGAReadLink{IV.7.3.19}{(7.3.19)}}, que sus fibras
formales sean geométricamente reducidas, o también que su completado
$\widehat A$ sea reducido y que los cuerpos componentes de su anillo total
de fracciones sean extensiones separables del cuerpo de fracciones
$K$ de $A$. En virtud de {\EGAReadLink{IV.7.6.4}{(7.6.4)}} y
{\EGAReadLink{IV.7.7.2}{(7.7.2)}}, esto equivale asimismo a
decir que $A$ es universalmente japonés. Cuando $A$ es un anillo de
valoración discreta, lo mismo ocurre con $\widehat A$; decir que $A$ es
excelente significa entonces que el cuerpo de fracciones de $\widehat A$
es una extensión separable de $K$, condición siempre satisfecha si
$K$ es de característica 0.

\paragraph{(Err\textsubscript{IV}, 28).}
Después de (IV, 8.10.3), después de la línea 14 de la p.~37, insertar:

\begin{remark}[8.10.3.1]
\label{IV.8.10.3.1}
En {\EGAReadLink{IV.8.10.3}{(8.10.3)}} no puede suprimirse
la hipótesis de que los $\overline{f_\alpha(X_\alpha)}$ sean
constructibles. Tomemos, por ejemplo, $S=\operatorname{Spec}(A)$, donde $A$
es un anillo tal que $S$ sea compacto, totalmente disconexo y sin
punto aislado, siendo maximal todo ideal primo $\mathfrak p$ de $A$ y
$A_\mathfrak p$ un cuerpo (Bourbaki, \emph{Alg. comm.}, cap.~II,
\S4, ejerc.~16 y 17). Sea $s\in S$, y tomemos como
$Y_\alpha=S_\alpha$ los entornos abiertos afines de $s$ en $S$; se
toma $X_\alpha=\operatorname{Spec}(k(s))$ para todo $\alpha$,
siendo $f_\alpha$ y $f$ los morfismos canónicos; está claro que
$f$, que es la identidad, es dominante, pero ninguno de los $f_\alpha$ lo
es.
\end{remark}

\paragraph{(Err\textsubscript{IV}, 29).}
En (IV, 8.12.8), línea 7 de la p.~47, sustituir « tel qu'il existe un
$\sh{O}_Y$-Module ample » por « et quasi-séparé ». Línea 12 de la
p.~47, sustituir « L'hypothèse » por « Il résulte de
{\EGAReadLink{IV.8.12.3}{(8.12.3)}} que l'existence de la
factorisation de l'énoncé est une propriété locale sur $Y$; on peut
donc supposer $Y$ affine, ce qui ».

\paragraph{(Err\textsubscript{IV}, 30).}
Después de (IV, 8.12.11), añadir al final de la p.~48:

\begin{corollary}[8.12.12]
\label{IV.8.12.12}
Sean $X$, $Y$ dos preesquemas íntegros,
$f:X\to Y$ un morfismo separado, birracional y localmente de tipo finito.
Sea $y\in Y$ un punto normal de $Y$ y supongamos que un punto
$x\in f^{-1}(y)$ esté aislado en $f^{-1}(y)$. Entonces existe un
entorno abierto $U$ de $y$ en $Y$ tal que la restricción
$f^{-1}(U)\to U$ de $f$ sea un isomorfismo.
\end{corollary}

\begin{proof}
Basta demostrar que existe un entorno abierto $U$ de $y$ en
$Y$ y una $U$-sección $g:U\to f^{-1}(U)$ de $f^{-1}(U)$ tal que
$g(y)=x$. En efecto, como $U$ contiene el punto genérico $\eta$ de
$Y$ y $f^{-1}(U)$ el punto genérico $\xi$ de $X$, se tiene
necesariamente $g(\eta)=\xi$ puesto que $f$ es birracional
{\EGAReadLink{IV.6.15.4}{(6.15.4)}}; como $g$ es una
inmersión, $g(U)$ es localmente cerrado en $f^{-1}(U)$, luego abierto
en su adherencia en $f^{-1}(U)$, y como esta última es igual
a $f^{-1}(U)$ por lo anterior, $g(U)$ es abierto en
$f^{-1}(U)$; como el subpreesquema de $f^{-1}(U)$ asociado a $g$ es
isomorfo a $U$, luego reducido, y $f^{-1}(U)$ es reducido, se deduce
de (I, 5.2.1) que $g$ es una inmersión abierta. Por otra parte, puesto que
$f$ es separado, $g$ es una inmersión cerrada (I, 5.4.6), luego $g(U)$
es un conjunto a la vez abierto y cerrado en $f^{-1}(U)$. Pero, como
$f^{-1}(U)$ es íntegro, esto implica $g(U)=f^{-1}(U)$, de donde resulta la
conclusión.

Consideremos entonces dos casos:

\textup{I)} Supongamos primero que $Y$ es normal, y sea $V$ el conjunto de los
puntos de $X$ en los cuales $f$ es cuasifinito; es un conjunto abierto
en $X$ {\EGAReadLink{IV.13.1.4}{(13.1.4)}}.\footnote{El
lector comprobará que ni {\EGAReadLink{IV.8.12.12}{(8.12.12)}}
ni ninguna de sus consecuencias se utiliza en la demostración de
{\EGAReadLink{IV.13.1.4}{(13.1.4)}}.} Se deduce entonces de
{\EGAReadLink{IV.8.12.10}{(8.12.10)}}

\oldpage[IV]{353}
que la restricción $f|V:V\to Y$ es una inmersión abierta; existe, por tanto,
un entorno abierto $U$ de $y\in Y$ tal que $f|V$ sea un isomorfismo
de $V$ sobre el subpreesquema inducido por $Y$ sobre $U$; el isomorfismo
recíproco $g$ es entonces una $U$-sección de $f^{-1}(U)$ tal que
$g(y)=x$.

\textup{II)} Como la obtención de una $U$-sección $g$ tal que $g(y)=x$ es
un problema local sobre $X$ y sobre $Y$, podemos suponer que $X$ e $Y$
son afines y que $f$ es de tipo finito; por tanto, $X$ es un
subpreesquema cerrado de un $Y$-preesquema $Z$ de presentación finita sobre
$Y$ (por ejemplo, de la forma $Y[T_1,\ldots,T_n]$). Pongamos
\[
 Y'=\operatorname{Spec}(\sh{O}_{Y,y}),\qquad
 X'=X_{(Y')},\qquad Z'=Z_{(Y')},
\]
de modo que $X'$ es un subpreesquema cerrado de $Z'$. Además, si
$f'=f_{(Y')}$, se tiene $f'^{-1}(y)=f^{-1}(y)$ (I, 3.6.5) y $x$ está, por tanto,
aislado en $f'^{-1}(y)$. Puesto que $Y'$ es normal, puede aplicarse el
resultado de I) a $f'$ y, puesto que $Y'$ es el único entorno de $y$ en
$Y'$, existe una $Y'$-sección $g'$ de $X'$ tal que $g'(y)=x$, que
evidentemente puede considerarse también como una $Y'$-sección de $Z'$.
Puesto que el morfismo estructural $h:Z\to Y$ es de presentación finita,
puede aplicarse ({\EGAReadLink{IV.8.8.2}{8.8.2}}, (i) y
(ii)) siguiendo el método de
({\EGAReadLink{IV.8.1.2}{8.1.2}}, a)), de donde se concluye
la existencia de un entorno abierto $U$ de $y$ en $Y$ y de una
$U$-sección $g$ de $h^{-1}(U)$ tal que $g(y)=x$. Queda por ver
que, de hecho, $g$ es una $U$-sección de $f^{-1}(U)$. Pero, puesto que
$f^{-1}(U)$ es un subpreesquema cerrado de $h^{-1}(U)$ y $U$ es
reducido, basta para ello (I, 5.2.2) comprobar que
$g(U)\subset f^{-1}(U)$. Ahora bien, puesto que $U$ es íntegro y $f^{-1}(U)$
es cerrado en $h^{-1}(U)$, basta demostrar que, si $\eta$ es el punto
genérico de $Y$ (y de $U$), se tiene $g(\eta)\in f^{-1}(U)$. Esto
resulta de que la restricción de $g$ a $Y'$ es una $Y'$-sección de
$X'\subset f^{-1}(U)$ y de que $\eta\in Y'$.
\end{proof}

\paragraph{(Err\textsubscript{IV}, 31).}
En ({\EGAReadLink{IV.10.4}{IV, 10.4}}), añadir, después de la
proposición {\EGAReadLink{IV.10.4.11}{(10.4.11)}}:

\begin{corollary}[10.4.12]
\label{IV.10.4.12}
Sean $S$ un preesquema, $X$ un $S$-preesquema de presentación finita y
$f$ un $S$-endomorfismo radicial de $X$.
\begin{enumerate}[label=\textup{(\roman*)}]
\item El endomorfismo $f$ es biyectivo y finito (por tanto, un homeomorfismo
universal {\EGAReadLink{IV.2.4.5}{(2.4.5)}}).
\item Supongamos además que $X$ sea reducido, que los anillos locales de
$S$ sean noetherianos y universalmente japoneses, y que $f$ sea
birracional (condición automáticamente satisfecha si, para todo punto
maximal $x$ de $X$, $k(x)$ es un cuerpo perfecto). Entonces $f$ es un
$S$-automorfismo.
\end{enumerate}
\end{corollary}

\begin{proof}
\textup{(i)} Como la cuestión es local sobre $S$, podemos limitarnos al caso
en que $S=\operatorname{Spec}(A)$ es afín. Existe entonces una
sub-$\mathbb Z$-álgebra de tipo finito $A_0$ de $A$ tal que, poniendo
$S_0=\operatorname{Spec}(A_0)$, se tenga
$X=X_0\times_{S_0}S$, donde $X_0$ es un preesquema de tipo finito sobre $S_0$,
y que $f$ se deduzca de un $S_0$-endomorfismo $f_0$ de $X_0$ mediante el
cambio de base $S\to S_0$
{\EGAReadLink{IV.8.9.1}{(8.9.1)}}; además, puede suponerse
$f_0$ radicial ({\EGAReadLink{IV.8.10.5}{8.10.5}}, (vii)).
Bastará demostrar que $f_0$ es biyectivo y finito, y podemos, por tanto,
suponer en lo sucesivo que $A$ es una $\mathbb Z$-álgebra de tipo finito, lo
cual implica que $S$ y $X$ son preesquemas excelentes
({\EGAReadLink{IV.7.8.3}{(7.8.3)}} y
{\EGAReadLink{IV.7.8.6}{(7.8.6)}}) y, en particular, que $X$
es un preesquema universalmente japonés
({\EGAReadLink{IV.5.11.4}{(5.11.4)}} y
({\EGAReadLink{IV.7.8.3}{7.8.3}}, (vi))). Además, puede
suponerse $X$ reducido; en efecto,
$f_{\mathrm{red}}:X_{\mathrm{red}}\to X_{\mathrm{red}}$ es un
$S_{\mathrm{red}}$-endomorfismo radicial (I, 5.1.6); si se demuestra que
$f_{\mathrm{red}}$ es finito, lo mismo ocurrirá con $f$: en efecto,
puesto que $f$ es radicial, es separado
{\EGAReadLink{IV.1.8.7.1}{(1.8.7.1)}}; como
$f_{\mathrm{red}}$ es propio, lo mismo ocurre con $f$ (II, 5.4.6);
pero, puesto que $f$ es evidentemente cuasifinito, es finito (III, 4.4.2).

\oldpage[IV]{354}
Sea entonces $X'$ la normalización de $X$ (II, 6.3.8); si $X_i$
($1\leq i\leq n$) son los subpreesquemas reducidos de $X$ cuyos
espacios subyacentes son las componentes irreducibles de $X$, $X'$ es
la suma de las normalizaciones $X'_i$ de los preesquemas íntegros $X_i$;
denotaremos por $x_i$ (resp. $x'_i$) el punto genérico de $X_i$
(resp. $X'_i$). El morfismo canónico $p:X'\to X$ es finito puesto que
$X$ es universalmente japonés y se tiene
$p^{-1}(x_i)=\{x'_i\}$, siendo biyectivo el homomorfismo
$k(x_i)\to k(x'_i)$.

Sea $f=(\psi,\theta)$; ya se sabe
{\EGAReadLink{IV.10.4.11}{(10.4.11)}} que $\psi$ es una
biyección continua del espacio $X$ sobre sí mismo; si se pone
$y_i=\psi(x_i)$, $y_i$ es también un punto maximal de $X$, pues si $y'_i$
fuese una generalización de $y_i$ distinta de $y_i$, $x_i$ sería una
generalización $\neq x_i$ de $\psi^{-1}(y'_i)$. Puesto que el conjunto $M$
de los puntos maximales de $X$ es finito, la restricción $\psi|M$ es una
biyección de $M$ sobre sí mismo, y puede escribirse
$\psi(x_i)=x_{\sigma(i)}$, donde $\sigma$ es una permutación de
$[1,n]$. Entonces, para todo $i$, la restricción $f|X_i$ de $f$ se
factoriza como
\[
 X_i\xrightarrow{f_i}X_{\sigma(i)}\longrightarrow X
\]
(I, 5.2.2), donde la segunda flecha es la inyección canónica. Como los
$X_i$ son íntegros, existe (II, 6.3.9) un único morfismo
$f'_i:X'_i\to X'_{\sigma(i)}$ que hace conmutativo el diagrama
\[
\begin{array}{ccc}
X'_i&\xrightarrow{f'_i}&X'_{\sigma(i)}\\
\downarrow p_i&&\downarrow p_{\sigma(i)}\\
X_i&\xrightarrow{f_i}&X_{\sigma(i)}
\end{array}
\]
y si $f'$ coincide con $f'_i$ en cada uno de los $X'_i$, se tiene, por tanto, un
diagrama conmutativo
\phantomsection\label{IV.10.4.12.1}
\begin{equation}
\label{IV.10.4.12.1.errata}
\begin{array}{ccc}
X'&\xrightarrow{f'}&X'\\
\downarrow p&&\downarrow p\\
X&\xrightarrow{f}&X.
\end{array}
\tag{10.4.12.1}
\end{equation}

Por otra parte, como $f$ es separado, lo mismo ocurre con
$f\circ p=p\circ f'$, luego $f'$ es separado (I, 5.5.1); por último, puesto que
$p$ es finito y $f$ biyectivo, se deduce de
{\EGAReadLink{IV.10.4.12.1}{(10.4.12.1)}} que $f'$ es
cuasifinito. Puesto que los $X'_i$ son íntegros y normales, se deduce de
{\EGAReadLink{IV.8.12.11}{(8.12.11)}} que $f'_i$ se factoriza como
\[
 X'_i\xrightarrow{f''_i}X''_{\sigma(i)}
 \xrightarrow{u_{\sigma(i)}}X'_{\sigma(i)},
\]
donde $X''_{\sigma(i)}$ es la clausura integral de
$X'_{\sigma(i)}$ relativa a la extensión $k(x'_i)$ de
$k(x'_{\sigma(i)})$, y $f''_i$ una inmersión abierta. Pero, como
$k(x'_i)$ es una extensión radicial de $k(x'_{\sigma(i)})$,
$u_{\sigma(i)}$ es un morfismo radicial (Bourbaki,
\emph{Alg. comm.}, cap.~V, \S2, n\textsuperscript{o}~3, lema 4);
luego $f'$ es radicial, y se deduce entonces de
{\EGAReadLink{IV.10.4.11}{(10.4.11)}} que $f'$ es
sobreyectivo; se concluye (puesto que los $u_{\sigma(i)}$ son inyectivos)
que $f''_i$ es biyectivo, luego un isomorfismo; por último, se deduce que
el morfismo $f'$ es finito (puesto que lo es $u_{\sigma(i)}$,
siendo $X'$ universalmente japonés
{\EGAReadLink{IV.7.8.3}{(7.8.3)}}). Ahora bien, como
$p\circ f'=f\circ p$ es entonces finito, $p$ es sobreyectivo y $f$
separado, $f$ es propio (II, 5.4.3) y cuasifinito, luego finito
{\EGAReadLink{IV.8.11.1}{(8.11.1)}}.

\oldpage[IV]{355}
\textup{(ii)} Teniendo en cuenta
({\EGAReadLink{IV.8.10.5}{8.10.5}}, (i)), aplicado siguiendo el
método de ({\EGAReadLink{IV.8.1.2}{8.1.2}}, a)), podemos
limitarnos al caso en que $S$ sea un esquema local. Las reducciones previas
realizadas en la demostración de (i) son, por tanto, innecesarias aquí, y con las mismas
notaciones se ve que se tiene un diagrama conmutativo
{\EGAReadLink{IV.10.4.12.1}{(10.4.12.1)}} en el cual $X$ es
un $S$-preesquema de tipo finito, $p$ es finito, $f'$ es un isomorfismo y
el homomorfismo canónico $\sh{O}_X\to p_*(\sh{O}_{X'})$ es inyectivo.
Pero se demostrará en el cap.~VI, en la teoría de los cocientes de
preesquemas por relaciones de equivalencia, que, bajo estas condiciones, $f$
es él mismo un isomorfismo (el lector comprobará que no se utilizará
{\EGAReadLink{IV.10.4.12}{(10.4.12)}} en la teoría de las
relaciones de equivalencia del cap.~VI).
\end{proof}

\begin{corollary}[10.4.13]
\label{IV.10.4.13}
Sean $k$ un cuerpo y $X$ un preesquema algebraico sobre $k$. Entonces todo
$k$-endomorfismo radicial $f$ de $X$ es finito y biyectivo, luego un
homeomorfismo universal. Si además $X$ es reducido y $f$ birracional
(por ejemplo, si $k$ es de característica 0), $f$ es un
$k$-automorfismo.
\end{corollary}

\begin{remark}[10.4.14]
\label{IV.10.4.14}
\begin{enumerate}[label=\textup{(\roman*)}]
\item En {\EGAReadLink{IV.10.4.12}{(10.4.12)}} y
{\EGAReadLink{IV.10.4.13}{(10.4.13)}} no puede suprimirse
la hipótesis de que $X$ es reducido para la segunda conclusión; así lo
muestra el ejemplo del $B$-endomorfismo correspondiente al aumento
$D_B(L)\to B$ del anillo $D_B(L)=A$ definido en ($0$, 18.2.3), donde
$B$ se identifica con un subanillo de $A$.

\item Sea $X$ un preesquema noetheriano. Se ignora si todo
endomorfismo radicial $f$ de $X$ es necesariamente sobreyectivo y finito.
Si se supone que $f$ es sobreyectivo y, además, $X$ es un preesquema
japonés, el razonamiento de
{\EGAReadLink{IV.10.4.12}{(10.4.12)}} muestra que $f$ es finito;
si además $X$ es reducido y $f$ es birracional, por ejemplo si,
para todo punto maximal $x\in X$, $k(x)$ es perfecto, el mismo
razonamiento muestra que $f$ es un automorfismo.

\item Pueden darse ejemplos de endomorfismos radiciales no
sobreyectivos de esquemas afines íntegros no noetherianos: si $k$ es un
anillo, $A=k[(T_\lambda)_{\lambda\in L}]$ un anillo de polinomios donde $L$
es infinito numerable, y $\mathfrak J$ el ideal engendrado en $A$ por
una subfamilia $(T_\lambda)_{\lambda\in H}$, donde $L-H$ es infinito y
$H\neq\varnothing$, se observa que $A/\mathfrak J$ es isomorfo a
$k[(T_\lambda)_{\lambda\in L-H}]$, luego a $A$; el endomorfismo de
$\operatorname{Spec}(A)$ correspondiente al endomorfismo de anillos
compuesto
\[
 A\xrightarrow{p}A/\mathfrak J\xrightarrow{u}A
\]
donde $p$ es la proyección canónica y $u$ un isomorfismo, responde a la
cuestión, siendo su imagen el subpreesquema cerrado de
$\operatorname{Spec}(A)$ definido por $\mathfrak J$.
\end{enumerate}
\end{remark}

\paragraph{(Err\textsubscript{IV}, 32).}
En (IV, 11.3.14), línea 7 de la p.~141, antes de « Soient », insertar
(i). Después de la línea 9 de la p.~141, añadir:

\smallskip
\noindent
\textup{(ii)} Recíprocamente, sea $f:X\to Y$ un morfismo plano. Si $X$
es unirrama (resp. geométricamente unirrama) en un punto $x$, $Y$
es unirrama (resp. geométricamente unirrama) en el punto $y=f(x)$.

Al comienzo de la línea 10 de la p.~141, insertar (i). Después de la línea 20
de la p.~141, añadir:

\smallskip
\noindent
\textup{(ii)} Observemos que, si $Y_1=Y_{\mathrm{red}}$,
$X_1=X\times_Y Y_1$ tiene como espacio subyacente el de $X$ y el
morfismo $f_1:X_1\to Y_1$ deducido de $f$ por cambio de base es
plano. Podemos, por tanto, limitarnos al caso en que $Y=\operatorname{Spec}(A)$,
siendo $A$ un anillo local reducido y $y$ el punto cerrado de $Y$. Pongamos
$B=\sh{O}_{X,x}$; por hipótesis $B_{\mathrm{red}}$ es unirrama, luego
$\operatorname{Spec}(B)$ tiene un único punto maximal, y por planitud
{\EGAReadLink{IV.2.3.4}{(2.3.4)}} se deduce que lo mismo ocurre

\oldpage[IV]{356}
con $\operatorname{Spec}(A)$; puesto que $A$ es reducido, es, por tanto,
íntegro. Sean $K$ el cuerpo de fracciones de $A$ y $A'$ la clausura
integral de $A$; por planitud, $B\otimes_AA'$ se identifica con un
subanillo de $B\otimes_AK$ que contiene $B$, y como los elementos no
nulos de $A$ son $B$-regulares por planitud, $B\otimes_AK$ se identifica
con un subanillo del anillo total de fracciones $R$ de $B$. Para
demostrar que $A$ es unirrama, basta ver que $A'$ no puede tener
dos ideales maximales distintos (necesariamente sobre $y$
(Bourbaki, \emph{Alg. comm.}, cap.~V, \S2,
n\textsuperscript{o}~1, prop.~1)); si así fuese,
$B'=B\otimes_AA'$ tendría dos ideales primos distintos sobre
$x$ (I, 3.4.7), y basta ver que esto no es posible.

Sea $\mathfrak N$ el nilradical de $B$; si $S$ es el conjunto de los
elementos regulares de $B$, se tiene $R=S^{-1}B$, y
$S^{-1}\mathfrak N=R\mathfrak N$ es el nilradical de $R$, luego
\[
 R/R\mathfrak N=S^{-1}B/S^{-1}\mathfrak N=S^{-1}(B/\mathfrak N)
\]
está contenido en el cuerpo de fracciones $L$ del anillo local íntegro
$B/\mathfrak N=B_{\mathrm{red}}$; lo mismo ocurre, por tanto, con
\[
 B'_{\mathrm{red}}=B'/(B'\cap R\mathfrak N)
\]
que es entero sobre $B_{\mathrm{red}}$; por hipótesis,
$B'_{\mathrm{red}}$, al estar contenido en la clausura integral de
$B_{\mathrm{red}}$, es un anillo local, luego lo mismo ocurre con
$B'$, lo cual establece nuestra afirmación.

Si se supone además que $B$ es geométricamente unirrama, el cuerpo
residual de $B'_{\mathrm{red}}$, al estar contenido en el cuerpo residual de
la clausura integral de $B_{\mathrm{red}}$, es una extensión radicial
del cuerpo residual $k(x)$ de $B$. Observemos ahora que el cuerpo residual
de $B'_{\mathrm{red}}$ es también el de $B'=B\otimes_AA'$, luego es
igual a $k(x)\otimes_{k(y)}k'$, donde $k'$ es el cuerpo residual de $A'$;
se deduce entonces de {\EGAReadLink{IV.2.6.1}{(2.6.1)}} que $k'$
es una extensión radicial de $k(y)$, luego $A$ es también
geométricamente unirrama.

\paragraph{(Err\textsubscript{IV}, 33).}
En ({\EGAReadLink{IV.12.1.2}{IV, 12.1.2}}), el ejemplo dado
en (i) no es correcto, pues el morfismo considerado es abierto. Se
construye del modo siguiente un ejemplo de un morfismo plano $f:X\to Y$ cuya restricción
a una componente irreducible de $X$ no es un morfismo abierto.
Sea $Z=\operatorname{Spec}(\mathbb C[U,V])$ el
plano afín sobre el cuerpo de los números complejos, y consideremos el
esquema $Y$ obtenido pegando dos puntos cerrados distintos $a$, $b$ de
$Z$; se tiene $Y=\operatorname{Spec}(B)$, donde $B$ es el subanillo de
$\mathbb C[U,V]$ formado por los polinomios que toman el mismo valor en los puntos
$a$ y $b$ (identificando el conjunto de los puntos cerrados de $Z$ con
$\mathbb C^2$); es inmediato que $B$ es un sub-$\mathbb C$-espacio
vectorial de $\mathbb C[U,V]$ de codimensión 1, luego $\mathbb C[U,V]$
es un $B$-módulo de tipo finito y, por consiguiente, el morfismo $Z\to Y$ es
finito, siendo la imagen inversa del punto $c$ de $Y$ obtenido pegando
$a$ y $b$ el conjunto $\{a,b\}$.

Consideremos entonces el esquema $X$ obtenido pegando dos ejemplares de
$Z$, sean $Z_1$, $Z_2$, de modo que, si $a_i$, $b_i$ ($i=1,2$) son
los puntos de $Z_i$ correspondientes a $a$ y $b$, $a_1$ se pegue con
$b_2$ y $b_1$ con $a_2$ (cf. ({\EGAReadLink{IV.6.5.5}{6.5.5}},
(ii))). Se ve fácilmente (\emph{loc. cit.}) que el morfismo $X\to Y$
así obtenido es plano y finito (e incluso étale) y que $Z_1$ y $Z_2$
se identifican con las dos componentes irreducibles de $X$. Sin embargo, si
$D$ es una recta en $Z_1$ que pasa por $a_1$ pero no por $b_1$,
la imagen $V$ en $Y$ del abierto $Z_1-D$ de $Z_1$ contiene la imagen de
$b_1$, es decir, el punto $c$, pero no contiene la imagen del punto
genérico de $D$, luego $V$ no es un entorno de $c$.

\paragraph{(Err\textsubscript{IV}, 34).}
En el título de (IV, 13.2), línea 5 desde abajo de la p.~190, sustituir
« équidimensionnels » por « localement équidimensionnels ». Del mismo modo,
en (IV, 13.2.2), línea 4 de la p.~191, sustituir dos veces
« équidimensionnel » por « localement équidimensionnel ».

\oldpage[IV]{357}
Línea 5 de la p.~191, añadir: \emph{Si además las fibras
$f^{-1}(f(x))$ de $f$ son equidimensionales ($0_{\mathrm{IV}}$,
14.1.3), se dice que $f$ es equidimensional.} Línea 6 desde abajo de la
p.~194, sustituir « équidimensionnelles » por « localement
équidimensionnelles ».

\paragraph{(Err\textsubscript{IV}, 35).}
En el título de (IV, 13.3), línea 4 desde abajo de la
p.~194\footnote{La lista de erratas impresa y el inglés mantenido dan
p.~94; el título de (IV, 13.3) se encuentra en la p.~194.}, sustituir
« équidimensionnels » por « localement équidimensionnels ». En
(IV, 13.3.2), línea 6 desde abajo de la p.~196, sustituir dos veces
« équidimensionnel » por « localement équidimensionnel ». Línea 5 desde abajo
de la p.~196, añadir: \emph{Si además las fibras $f^{-1}(f(x))$ de
$f$ son equidimensionales ($0_{\mathrm{IV}}$, 14.1.3), se dice que
$f$ es equidimensional.} En (IV, 13.3.7), línea 14 de la p.~198,
sustituir « équidimensionnel » por « localement équidimensionnel ».
En ({\EGAReadLink{IV.13.3.9}{IV, 13.3.9}}), línea 3 desde
abajo de la p.~198, sustituir « équidimensionalité » por
« équidimensionalité locale ». Líneas 10, 12, 20 y 24 de la p.~199,
sustituir « équidimensionnel » por « localement équidimensionnel ».

\paragraph{(Err\textsubscript{IV}, 36).}
En la introducción del \S14, línea 6 desde abajo de la p.~199, sustituir
« équidimensionnel » por « localement équidimensionnel ».

\paragraph{(Err\textsubscript{IV}, 37).}
En (IV, 14.4.1), línea 5 desde abajo de la p.~209, sustituir « localement
de type fini » por « localement de présentation finie ». El enunciado
actual de ({\EGAReadLink{IV.14.4.1}{IV, 14.4.1}}) es inexacto,
incluso cuando $f:X\to Y$ es una inmersión cerrada, $X$ está reducido a
un solo punto $y$ (cerrado) de $Y$ y $\sh{O}_{Y,y}$ es un cuerpo. En
efecto, existen anillos $A$ tales que $Y=\operatorname{Spec}(A)$ es
compacto, no tiene puntos aislados y es totalmente disconexo, todo ideal primo
$\mathfrak p$ de $A$ es maximal y $A_\mathfrak p$ es un cuerpo
(Bourbaki, \emph{Alg. comm.}, cap.~II, \S4, ejerc.~16 y 17); es
claro que la inmersión cerrada $\{y\}\to Y$ no es un morfismo
abierto, pero cumple las condiciones b) y b') de
({\EGAReadLink{IV.14.4.1}{IV, 14.4.1}}).

\paragraph{(Err\textsubscript{IV}, 38).}
En (IV, 14.4.1.3), línea 14 de la p.~211, sustituir « localement
quasi-fini et dominant » por « localement quasi-fini, localement de
présentation finie et dominant ».

\paragraph{(Err\textsubscript{IV}, 39).}
En (IV, 14.4.2), línea 10 desde abajo de la p.~211, sustituir « localement
de type fini » por « localement de présentation finie ».

\paragraph{(Err\textsubscript{IV}, 40).}
En (IV, 14.4.3), línea 18 de la p.~212, sustituir « localement de
type fini » por « localement de présentation finie ».

\paragraph{(Err\textsubscript{IV}, 41).}
En ({\EGAReadLink{IV.14.4.4}{IV, 14.4.4}}), líneas 15 y 16 desde
abajo de la p.~212, sustituir « localement de type fini » por
« localement de présentation finie ». Líneas 10 y 11 desde abajo, sustituir
« équidimensionnel » por « localement équidimensionnel ».

\paragraph{(Err\textsubscript{IV}, 42).}
En (IV, 14.4.8), después de la línea 1 de la p.~214, añadir:
\begin{quote}
\emph{Si además el conjunto de las componentes irreducibles de $Y$ es
localmente finito, estas dos condiciones son equivalentes a las siguientes:}
\end{quote}

\paragraph{(Err\textsubscript{IV}, 43).}
En (IV, 14.4.10), línea 15 de la p.~215, sustituir « quelconque » por
« dont l'ensemble des composantes irréductibles est localement fini ».
Línea 16 de la p.~215, sustituir « de type fini » por « de présentation
finie ».

\paragraph{(Err\textsubscript{IV}, 44).}
En (IV, 15.5.6.1), línea 15 desde abajo de la p.~234, después de la primera
frase, añadir:

\begin{quote}
Podemos limitarnos al caso en que toda parte cerrada irreducible de $X$
admite un único punto genérico. En efecto, en caso contrario consideramos
el espacio $Y=X/R$,

\oldpage[IV]{358}
cociente de $X$ por la relación de equivalencia
$\overline{\{x\}}=\overline{\{x'\}}$. De esta definición se sigue
inmediatamente que todo abierto de $X$ es saturado para $R$; por tanto, la
topología de $X$ es la imagen recíproca de la de $Y$ por la aplicación
canónica $\pi:X\to X/R$. Un abierto denso de $X$ es de la forma
$\pi^{-1}(U)$, donde $U$ es un abierto denso de $Y$; por consiguiente, una
parte cerrada rara $Z$ de $X$ es de la forma $\pi^{-1}(S)$, donde
$S=\pi(Z)$ es cerrada y rara en $Y$, y decir que $X-Z$ es conexo equivale
a decir que $Y-S$ es conexo. Del mismo modo, $\pi(T_x)$ es el conjunto de
las generalizaciones de $\pi(x)$, y decir que $T_x-\{x\}$ es conexo
equivale a decir que $\pi(T_x)-\{\pi(x)\}$ lo es. Finalmente, por
construcción, $Y$ es un espacio de Kolmogorov localmente noetheriano, por lo
que toda parte cerrada irreducible admite un único punto genérico.
Se ve entonces que demostrar el lema para $X$ equivale a demostrarlo para
$Y$, y podemos suponer que toda parte cerrada irreducible de $X$ tiene un
único punto genérico.
\end{quote}

\paragraph{(Err\textsubscript{IV}, 45).}
En (IV, 15.6.1), línea 18 desde abajo de la p.~236, suprimir « localement
noethérien » (hipótesis no utilizada en la demostración).

\paragraph{(Err\textsubscript{IV}, 46).}
En (IV, 15.6.3), línea 11 de la p.~237, en lugar de « supposons que
pour tout $y\in Y$ », leer: « supposons que $f$ soit localement de
présentation finie et que pour tout $y\in Y$ ». Sustituir la
demostración (líneas 14 a 24 de la p.~237) por la siguiente:
\begin{quote}
En virtud de {\EGAReadLink{IV.1.10.3}{(1.10.3)}}, basta demostrar
que para todo $x\in X_y^0$ y toda generalización $y'$ de $y$, existe una
generalización $x'$ de $x$ tal que $f(x')=y'$. Ahora bien, basta tomar
como $x'$ el punto genérico de $X_{y'}^0$, en virtud de
{\EGAReadLink{IV.15.6.2}{(15.6.2)}}.
\end{quote}

\paragraph{(Err\textsubscript{IV}, 47).}
En (IV, 15.6.7), línea 2 de la p.~241, después de « noethérien », añadir:
« ou si la condition $\beta$) de
({\EGAReadLink{IV.15.6.6}{15.6.6}}, (iii)) est vérifiée ».
Después de la línea 9 de la p.~241, añadir:
\begin{quote}
Cuando se cumple la condición $\beta$) de
({\EGAReadLink{IV.15.6.6}{15.6.6}}, (iii)), el hecho de que
a) implica b') resulta de
({\EGAReadLink{IV.15.6.6}{15.6.6}}, (ii)), y el hecho de que
b') implica a) resulta de
({\EGAReadLink{IV.15.6.6}{15.6.6}}, (iii)).
\end{quote}

\paragraph{(Err\textsubscript{IV}, 48).}
En (IV, 20.1.5), línea 13 de la p.~228, (IV, 20.1.6), línea 17 de la
p.~228, y (IV, 20.1.7), línea 13 desde abajo de la p.~228, suprimir
« strictement ».

\paragraph{(Err\textsubscript{IV}, 49).}
En (IV, 20.2.13), añadir después de la línea 6 de la p.~236: En
particular, si $X$ es reducido y solo tiene un número finito de
componentes irreducibles, la noción de $\sh{O}_X$-módulo sin torsión
definida en ({\EGAReadLink{IV.20.1.5}{IV, 20.1.5}})
coincide con la de (I, 7.4.1).

\paragraph{(Err\textsubscript{IV}, 50).}
En (IV, 20.2.14), líneas 18 y 25 de la p.~236, suprimir
« strictement ».

\paragraph{(Err\textsubscript{IV}, 51).}
En (IV, 20.6.2), línea 13 desde abajo de la p.~252, suprimir
« strictement ».

\paragraph{(Err\textsubscript{IV}, 52).}
En (IV, 20.6.4), líneas 23 y 22 de la p.~253, suprimir
« strictement ».

\paragraph{(Err\textsubscript{IV}, 53).}
En (IV, 21.4), después de la línea 23 de la p.~267, añadir:

\begingroup
\def\EGAIVSuppressErratumLabels{1}
% EGA IV, Errata y Addenda, lista 3, n.º 53.
% Este componente se inserta después de 21.4.8 en el texto de lectura corregido
% y se repite literalmente bajo Err_IV, 53 en el apéndice de erratas impreso.

La proposición siguiente y sus corolarios nos fueron señalados por
M.~Raynaud.

\begin{proposition}[21.4.9]
\ifdefined\EGAIVSuppressErratumLabels\else\label{IV.21.4.9}\fi
Sean $X$ e $Y$ dos preesquemas localmente noetherianos, y sea
$f:X\to Y$ un morfismo fielmente plano. Supongamos además que $f$ es
cuasicompacto o que $f$ es abierto. Un divisor $D'$ sobre $X$ es de la
forma $f^*(D)$, con $D$ un divisor sobre $Y$, si y solo si se cumple la
condición siguiente: para todo $y\in Y$ tal que
$\operatorname{prof}(\sh{O}_{Y,y})\leq 1$, pongamos
\[
 Y_1=\operatorname{Spec}(\sh{O}_{Y,y}),\qquad
 X_1=X\times_Y Y_1,
 \qquad f_1=f_{(Y_1)}:X_1\to Y_1;
\]
entonces la imagen recíproca $D_1'$ de $D'$ sobre $X_1$ es de la forma
$f_1^*(D_1)$ para algún $D_1\in\operatorname{Div}(Y_1)$.
\end{proposition}

\oldpage[IV]{359}
\begin{proof}
La necesidad de la condición resulta inmediatamente de la transitividad de las
imágenes recíprocas de divisores, pues todos los morfismos considerados son planos
{\EGAReadLink{IV.21.4.4}{(21.4.4)}}. Para todo abierto $U$ de $Y$, denotemos por $f_U$ la
restricción $f^{-1}(U)\to U$ de $f$. Por {\EGAReadLink{IV.21.4.8}{(21.4.8)}}, la aplicación
\[
 D_U\longmapsto f_U^*(D_U)
\]
de $\operatorname{Div}(U)$ en
$\operatorname{Div}(f^{-1}(U))$ es inyectiva. Así, si $U_1,U_2$ son
dos abiertos de $Y$ y
\[
 D'|f^{-1}(U_1)=f_{U_1}^*(D_1),\qquad
 D'|f^{-1}(U_2)=f_{U_2}^*(D_2),
\]
entonces necesariamente
$D_1|U_1\cap U_2=D_2|U_1\cap U_2$. Se sigue inmediatamente que existe un
mayor abierto $U\subset Y$ para el cual
\[
 D'|f^{-1}(U)=f_U^*(D_U)
 \quad\text{con}\quad D_U\in\operatorname{Div}(U).
\]
Debemos demostrar que $U=Y$. Supongamos lo contrario, y sea $y$ un punto
maximal de $Y\setminus U$. Encontraremos un entorno abierto $V$ de
$y$ en $Y$ tal que
  $D'|f^{-1}(V)=f_V^*(D_V)$ para algún
$D_V\in\operatorname{Div}(V)$, contradiciendo la maximalidad de $U$.

\textup{I) Caso en que $f$ es cuasicompacto.}
Utilizando la cuasicompacidad de $f$ y {\EGAReadLink{IV.20.3.8}{(20.3.8)}}, por el método de
({\EGAReadLink{IV.8.1.2}{8.1.2}}, a), nos reducimos a demostrar ---con las notaciones del
enunciado, pero sin imponer inicialmente una cota a
$\operatorname{prof}(\sh{O}_{Y,y})$--- que $D_1'$ es de la forma
$f_1^*(D_1)$. Si $\operatorname{prof}(\sh{O}_{Y,y})\leq 1$, esto es
precisamente la hipótesis. Supongamos, pues, que
$\operatorname{prof}(\sh{O}_{Y,y})\geq 2$. Como $y$ es un punto
maximal de $Y\setminus U$, el abierto $W=Y_1\cap U$ de $Y_1$ es
$Y_1\setminus\{y\}$. Podemos, por tanto, limitarnos al caso
$Y=Y_1$ y $U=Y\setminus\{y\}$.

Por hipótesis existe un divisor $D_U\in\operatorname{Div}(U)$ tal
que $f_U^*(D_U)=D'|f^{-1}(U)$. El divisor $D_U$ es la clase de
equivalencia de un par $(\sh{L}_U,s)$, donde $\sh{L}_U$ es un
$\sh{O}_U$-módulo invertible y $s$ una sección meromorfa regular de
$\sh{L}_U$ sobre $U$ {\EGAReadLink{IV.21.2.11}{(21.2.11)}}. Por tanto, en virtud de
{\EGAReadLink{IV.21.4.1}{(21.4.1)}}, $D'|f^{-1}(U)$ es la clase de equivalencia de
$(\sh{L}'_U,s')$, donde
\[
 \sh{L}'_U=f_U^*(\sh{L}_U),\qquad s'=s\circ f_U.
\]
Sean $i:U\to Y$ y $j:f^{-1}(U)\to X$ las inmersiones canónicas,
y pongamos
\[
 \sh{L}=i_*(\sh{L}_U),\qquad
 \sh{L}'=j_*(\sh{L}'_U).
\]
Como $f$ es plano, $f^*(\sh{L})$ es canónicamente isomorfo a
$\sh{L}'$ {\EGAReadLink{IV.2.3.1}{(2.3.1)}}. Además, la planitud y la desigualdad
$\operatorname{prof}(\sh{O}_{Y,y})\geq 2$ implican
$\operatorname{prof}(\sh{O}_{X,x})\geq 2$ para todo
$x\in f^{-1}(y)$ {\EGAReadLink{IV.6.3.1}{(6.3.1)}}.

El divisor $D'$ es la clase de equivalencia de un par
$(\sh{L}'',s'')$, donde $\sh{L}''$ es un $\sh{O}_X$-módulo invertible
y $s''$ una sección meromorfa regular de $\sh{L}''$ sobre $X$.
Por consiguiente, $\sh{L}'_U$ es isomorfo a
$\sh{L}''|f^{-1}(U)$, y las observaciones precedentes, junto con
{\EGAReadLink{IV.5.9.8}{(5.9.8)}} y {\EGAReadLink{IV.5.10.5}{(5.10.5)}}, muestran que $\sh{L}'$ es un
$\sh{O}_X$-módulo invertible isomorfo a $\sh{L}''$. Como $f$ es
fielmente plano y cuasicompacto, se sigue que $\sh{L}$ es un
$\sh{O}_Y$-módulo invertible {\EGAReadLink{IV.2.5.2}{(2.5.2)}} y, por tanto, isomorfo a $\sh{O}_Y$, pues
$Y$ es un esquema local.

Además, como $\operatorname{prof}(\sh{O}_{Y,y})\geq 1$, el abierto
$U=Y\setminus\{y\}$ es esquemáticamente denso en $Y$
({\EGAReadLink{IV.20.2.13}{20.2.13}}, (iv)). La sección meromorfa regular $s$ de $\sh{L}_U$
sobre $U$ se prolonga, por tanto, a una sección meromorfa regular $s_0$ de
$\sh{L}$ sobre $Y$ {\EGAReadLink{IV.20.2.11}{(20.2.11)}}. La clase de equivalencia de
$(\sh{L},s_0)$ es así un divisor $D$ sobre $Y$. Evidentemente, $f^*(D)$ y
$D'$ inducen el mismo divisor sobre
$f^{-1}(U)=X\setminus f^{-1}(y)$; y, como
$\operatorname{prof}(\sh{O}_{X,x})\geq 2$ para todo
$x\in f^{-1}(y)$, {\EGAReadLink{IV.21.1.8}{(21.1.8)}} da $D'=f^*(D)$. Esto prueba el caso I.

\textup{II) Caso en que $f$ es abierto.}
Se trata de encontrar un entorno abierto $V$ de $y$ en $Y$ tal
que $D'|f^{-1}(V)=f_V^*(D_V)$ para algún
$D_V\in\operatorname{Div}(V)$, por lo que podemos suponer primero que $Y$ es afín.
Como $f$ es abierto, las imágenes $f(W)$, cuando $W$ recorre los
abiertos afines de $X$, forman un recubrimiento abierto de $Y$. Como $Y$
es cuasicompacto, existe un abierto cuasicompacto $T\subset X$
tal que $f(T)=Y$, y $f|T$ es cuasicompacto (I, 1.1.1). Con mayor razón,
para todo abierto cuasicompacto $T_\alpha\supset T$ de $X$,
\oldpage[IV]{360}
se tiene
$f(T_\alpha)=Y$ y $f|T_\alpha$ es cuasicompacto.

Aplicando el caso I a $f|T_\alpha$, obtenemos un divisor $D_\alpha$ sobre
$Y$ tal que
\[
 (f|T_\alpha)^*(D_\alpha)=D'|T_\alpha.
\]
Por transitividad de las imágenes recíprocas de divisores por morfismos planos
{\EGAReadLink{IV.21.4.4}{(21.4.4)}}, también se tiene
$(f|T)^*(D_\alpha)=D'|T$. Por tanto, para cualesquiera $\alpha,\beta$,
\[
 (f|T)^*(D_\alpha)=(f|T)^*(D_\beta),
\]
y {\EGAReadLink{IV.21.4.8}{(21.4.8)}} da $D_\alpha=D_\beta$. Si $D$ denota su valor
común, entonces $f^*(D)|T_\alpha=D'|T_\alpha$ para todo $\alpha$, y
por consiguiente $f^*(D)=D'$.
\end{proof}

\begin{corollary}[21.4.10]
\ifdefined\EGAIVSuppressErratumLabels\else\label{IV.21.4.10}\fi
Supongamos las hipótesis generales de {\EGAReadLink{IV.21.4.9}{(21.4.9)}} y, además, que
$Y$ es normal y que, para todo $y\in Y$ con
$\dim(\sh{O}_{Y,y})=1$, la fibra $X_y=f^{-1}(y)$ no tiene ningún
ciclo primo asociado inmerso. Un divisor $D'$ sobre $X$ es de la forma $f^*(D)$,
con $D\in\operatorname{Div}(Y)$, si y solo si se cumplen las dos
condiciones siguientes:
\begin{enumerate}[label=\textup{\arabic*\textdegree}]
\item para todo punto maximal $y$ de $Y$, la imagen recíproca de $D'$ en
$\operatorname{Div}(X_y)$ es cero;
\item para todo $y\in Y$ tal que $\dim(\sh{O}_{Y,y})=1$, existe un
entero $n_y$ tal que, para todo punto maximal $x$ de $X_y$, la
imagen recíproca de $D'$ en
$\operatorname{Div}(\operatorname{Spec}(\sh{O}_{X,x}))$ es igual a
$\operatorname{div}(t_y^{n_y})$, donde $t_y$ es la imagen por
$\sh{O}_{Y,y}\to\sh{O}_{X,x}$ de un uniformizante del anillo de
valoración discreta $\sh{O}_{Y,y}$.
\end{enumerate}
\end{corollary}

\begin{proof}
Aplicamos el criterio {\EGAReadLink{IV.21.4.9}{(21.4.9)}}. En un punto maximal $y$ de $Y$,
$\operatorname{Div}(Y_1)=0$ porque $\sh{O}_{Y,y}$ es artiniano;
la condición 1\textdegree{} es, por tanto, equivalente al criterio de
{\EGAReadLink{IV.21.4.9}{(21.4.9)}} en tal punto. Queda por demostrar que las condiciones
1\textdegree{} y 2\textdegree{} implican asimismo el criterio en un
punto $y\in Y$ con $\operatorname{prof}(\sh{O}_{Y,y})=1$. Como $Y$
es normal, estos son precisamente los puntos para los cuales
$\dim(\sh{O}_{Y,y})=1$, equivalentemente, aquellos para los cuales
$\sh{O}_{Y,y}$ es un anillo de valoración discreta {\EGAReadLink{IV.5.8.6}{(5.8.6)}}.

Podemos, pues, limitarnos a $Y=\operatorname{Spec}(\sh{O}_{Y,y})$. Si $t$
es un uniformizante de $\sh{O}_{Y,y}$, la aplicación
$n\mapsto\operatorname{div}(t^n)$ es un isomorfismo de $\mathbb Z$
sobre $\operatorname{Div}(Y)$. La condición de {\EGAReadLink{IV.21.4.9}{(21.4.9)}} dice que
$D'$ es de la forma $\operatorname{div}(t^n\circ f)$ para algún
$n\in\mathbb Z$, e implica evidentemente la condición 2\textdegree{}.
Recíprocamente, supongamos que se cumple 2\textdegree{}, así como 1\textdegree{}, y
escribamos $n=n_y$. Queda por demostrar
$D'=\operatorname{div}(t^n\circ f)$. Si $\eta$ es el punto genérico de
$Y$, los gérmenes de estos dos divisores en todo punto de $X_\eta$ son
cero por la condición 1\textdegree{}. Sus gérmenes en todo punto maximal
de $X_y$ son iguales por la condición 2\textdegree{}. En un punto no maximal
$x\in X_y$, se tiene
$\operatorname{prof}(\sh{O}_{X_y,x})\geq 1$, porque $X_y$ no tiene
ciclo primo asociado inmerso y, por tanto,
$\operatorname{prof}(\sh{O}_{X,x})\geq 2$ por {\EGAReadLink{IV.6.3.1}{(6.3.1)}}. La
igualdad buscada resulta de {\EGAReadLink{IV.21.1.8}{(21.1.8)}}.
\end{proof}

\begin{corollary}[21.4.11]
\ifdefined\EGAIVSuppressErratumLabels\else\label{IV.21.4.11}\fi
Supongamos las hipótesis generales de {\EGAReadLink{IV.21.4.9}{(21.4.9)}} y, además, que
$Y$ es normal y que, para todo $y\in Y$ con
$\dim(\sh{O}_{Y,y})=1$, la fibra $X_y$ es íntegra. Un divisor $D'$ es
de la forma $f^*(D)$, con $D\in\operatorname{Div}(Y)$, si y solo si,
para todo punto maximal $y$ de $Y$, la imagen recíproca de $D'$ en
$\operatorname{Div}(X_y)$ es cero.
\end{corollary}

\begin{proof}
Es suficiente comprobar que la condición 2\textdegree{} de {\EGAReadLink{IV.21.4.10}{(21.4.10)}} es entonces
automática. Si $\dim(\sh{O}_{Y,y})=1$ y $X_y$ es íntegra, entonces en
el único punto maximal $x$ de $X_y$ se tiene
$\dim(\sh{O}_{X,x})=1$ por {\EGAReadLink{IV.6.1.3}{(6.1.3)}}. Además,
$\sh{O}_{Y,y}$ es regular y $\sh{O}_{X_y,x}$ es un cuerpo, por lo que
$\sh{O}_{X,x}$ es regular {\EGAReadLink{IV.6.5.1}{(6.5.1)}} y, por tanto, un anillo de valoración discreta.
Si $t$ es un uniformizante de $\sh{O}_{Y,y}$, su imagen $t_y$ en
$\sh{O}_{X,x}$ es un uniformizante de $\sh{O}_{X,x}$
($0$, 17.3.3). La imagen recíproca de $D'$ en
$\operatorname{Div}(\operatorname{Spec}(\sh{O}_{X,x}))$ es, por tanto,
\oldpage[IV]{361}
de la forma $\operatorname{div}(t_y^n)$, lo que es exactamente la condición
2\textdegree{} de {\EGAReadLink{IV.21.4.10}{(21.4.10)}}, puesto que $X_y$ solo tiene un punto maximal.
\end{proof}

\begin{remark}[21.4.12]
\ifdefined\EGAIVSuppressErratumLabels\else\label{IV.21.4.12}\fi
Si en {\EGAReadLink{IV.21.4.11}{(21.4.11)}} se supone solamente que los $X_y$ son reducidos, hay que
añadir a la condición del enunciado que las imágenes recíprocas de $D'$
en cada uno de los grupos
$\operatorname{Div}(\operatorname{Spec}(\sh{O}_{X,x}))$, cuando $x$
recorre los puntos maximales de $X_y$, sean todas de la forma
$\operatorname{div}((t_x')^n)$, donde $t_x'$ es un uniformizante de
$\sh{O}_{X,x}$ y el entero $n$ es el mismo para todos estos puntos.
\end{remark}

\begin{corollary}[21.4.13]
\ifdefined\EGAIVSuppressErratumLabels\else\label{IV.21.4.13}\fi
Supongamos las hipótesis de {\EGAReadLink{IV.21.4.11}{(21.4.11)}}. Para todo
$\sh{O}_X$-módulo invertible $\sh{L}$ tal que, para todo punto maximal $\eta$ de
$Y$, la imagen recíproca
\[
 \sh{L}_\eta
 =\sh{L}\otimes_{\sh{O}_X}\sh{O}_{X_\eta}
\]
sea isomorfa a $\sh{O}_{X_\eta}$, existen un
$\sh{O}_Y$-módulo invertible $\sh{M}$ y un isomorfismo
$f^*(\sh{M})\simeq\sh{L}$.
\end{corollary}

\begin{proof}
Basta encontrar una sección meromorfa regular $s$ de $\sh{L}$ sobre
$X$ tal que el divisor $D'$ correspondiente a $(\sh{L},s)$ tenga
$D'_x=0$ en todo punto $x$ de una fibra $X_\eta$ sobre un punto maximal
$\eta$ de $Y$. Aplicando entonces {\EGAReadLink{IV.21.4.11}{(21.4.11)}} a $D'$, se obtiene el resultado por
{\EGAReadLink{IV.21.4.2}{(21.4.2)}}.

Los puntos de $\operatorname{Ass}(\sh{O}_X)$ están sobre los puntos
maximales de $Y$ {\EGAReadLink{IV.3.3.1}{(3.3.1)}}. Como $Y$ es reducido y localmente noetheriano,
existe un abierto denso $V$ de $Y$ que es la unión de una
familia disjunta de abiertos afines, cada uno de los cuales contiene exactamente un
punto maximal de $Y$; y $f^{-1}(V)$ es esquemáticamente denso en $X$
{\EGAReadLink{IV.5.10.2}{(5.10.2)}}. Podemos, por tanto, limitarnos al caso en que $Y$ es
afín e íntegro.

Para el único punto maximal $\eta$ de $Y$, la fibra $X_\eta$ es entonces
íntegra, de modo que $\operatorname{Ass}(\sh{O}_X)$ consta de un solo punto
$z$, el punto genérico de $X$. Sea $s_\eta$ la sección de
$\sh{L}_\eta$ sobre $X_\eta$ que corresponde, mediante un isomorfismo,
a la sección unidad de $\sh{O}_{X_\eta}$. Sea $U$ un entorno abierto
afín en $X$ de un punto $x\in X_\eta$, y sea $s$ una sección
de $\sh{L}$ sobre $U$ con el mismo germen en $x$ que $s_\eta$. Como $U$
es esquemáticamente denso en $X$ ({\EGAReadLink{IV.20.2.13}{20.2.13}}, (iv)), $s$ se identifica con una
sección meromorfa regular unívocamente determinada de $\sh{L}$ sobre $X$.
Resulta claramente de {\EGAReadLink{IV.20.2.11}{(20.2.11)}} y {\EGAReadLink{IV.20.3.5}{(20.3.5)}} que esta sección induce sobre
$X_\eta$ una sección meromorfa regular de $\sh{L}_\eta$ igual a
$s_\eta$ y, por tanto, tiene la propiedad requerida.
\end{proof}

\endgroup

\bigskip
\begin{center}
\small
1967. --- Imprimerie des Presses Universitaires de France. --- Vendôme
(France)\\
ÉDIT. N\textsuperscript{o} 29 470\qquad
IMPRIMÉ EN FRANCE\qquad
IMP. N\textsuperscript{o} 20 298
\end{center}
